The recession velocity is calculated by either the classical relation,
$v_c = z\times c$, or the relativistic relation,
\[
v_r = \frac{(1+z)^2 - 1}{(1+z)^2 + 1}\,c.
\]
The calculations for the three galaxies are summarized in the following
Table:

\begin{center}
\begin{tabular}{|c|c|c|c|c|}
\hline
Galaxy & $v_c$ (km/s) & $V_a$ (km/s) & $v_r$ (km/s) & $v_r/c \times 100$\\
\hline
3C279 & 160800 & 128750 & 121390 & 40\%\\
\hline
3C245 & 308700 & 212260 & 182740 & 61\%\\
\hline
4C41.17 & 1140000 & 470580 & 275040 & 92\%\\
\hline
\end{tabular}
\end{center}

\begin{description}
\item[(1)] Table 1, Columns 2,3,4 \hfill(5 Points)
\item[(2)] Column 5 \hfill(3 Points)
\item[(3)] If the student answers ``Classical'' \hfill(0 Points)\\
  If the student answers ``Special Relativity'' \hfill(1 Point)\\
  If the student answers ``Approximate formula'' \hfill(2 Points)
\end{description}
